\section{A conditional square-root-scale deficit for partial Hadamard matrices}
\label{sec:conditional-half-frame}

Let $r(n)$ denote the largest number of orthogonal sign rows of length
$n$, and let $\tau(u)$ be the least exponent $t$ for which a Hadamard
matrix of order $2^tu$ exists, for odd $u$. The established multibase
input is unconditional:
\[
 \tau(u)=O\!\left(\frac{\log u}{\log\log u}\right).
\]
We combine it with a \emph{conditional} prime-distribution estimate.
Throughout, logarithms are natural. No assertion that GRH has been
proved or verified is made.

The concatenation strategy is historical: de Launey--Gordon,
\emph{A Comment on the Hadamard Conjecture}, JCTA 95 (2001),
180--184, Lemma 2.1 and its proof,
\url{https://doi.org/10.1006/jcta.2000.3155}, use a Paley block and
an asymptotic-existence block in a short progression. The conclusion
below tracks the new sublogarithmic multiplier bound rather than
replacing it by an arbitrary fixed power. No novelty claim is made
pending a dedicated literature comparison.

\begin{theorem}[Conditional quantitative half-frame]
Assume GRH for primitive Dirichlet $L$-functions whose conductors are
powers of two, including conductor one. There are positive constants
$B,K$ such that every sufficiently large $n\equiv0\pmod4$ satisfies
\[
 \boxed{r(n)\ge\frac n2-
 B\sqrt n\exp\!\left(K\frac{\log n}{\log\log n}\right).}
\]
In particular the conditional deficit is $n^{1/2+o(1)}$ with a
specified subpower envelope, not merely $O_\epsilon(n^{1/2+\epsilon})$.
The theorem remains conditional even though its $\tau$ input is
unconditional.
\end{theorem}

\subsection{A uniform multiplier buffer removes circularity}
\begin{lemma}[Uniform power-of-two buffer]
There is an integer-valued function $T(n)\ge2$ such that
\[
 T(n)=O\!\left(\frac{\log n}{\log\log n}\right),
 \qquad \tau(u)\le T(n)\quad
 \text{for every positive odd }u\le n.
\]
For $Q(n)=2^{T(n)}$, every integer $1\le k\le n$ has a Hadamard
matrix of order $Q(n)k$.
\end{lemma}
\begin{proof}
Choose $A,u_0$ so $\tau(u)\le A\log u/\log\log u$ for odd
$u\ge u_0$, increasing $u_0$ beyond $e^e$. The function
$\log x/\log\log x$ is increasing there. The standard all-odd
multiplier construction ensures $\tau(u)<\infty$ also for the finitely
many smaller odd parts (for example, the binary-length complementary
pair construction supplies them). A fixed additive integer buffer
covers these smaller odd $u$, so
$T(n)=\lceil A\log n/\log\log n\rceil+B_0$ works for sufficiently
large $n$, with $B_0\ge2$. Write $k=2^vu$ with $u$ odd. Then
$\tau(u)\le T(n)$, and Sylvester doubling raises order
$2^{\tau(u)}u$ to $2^{T(n)+v}u=Q(n)k$.
\end{proof}

In particular $Q(n)=n^{o(1)}$. The buffer is chosen from $n$ before
the progression element or $k$ is selected. No bound on $\tau(k)$
depending on an as-yet-unknown construction is being assumed.

\subsection{The precise conditional number-theory input}
Put
\[
 \psi(x;Q,a)=\sum_{\substack{m\le x\\m\equiv a\pmod Q}}\Lambda(m).
\]
Under Dirichlet GRH, uniformly for reduced residues and, in particular,
for $Q\le\sqrt x$,
\[
 \psi(x;Q,a)=\frac{x}{\varphi(Q)}+
 O\bigl(\sqrt x(\log x)^2\bigr),
\]
with an absolute implied constant. A primary author source is
Montgomery--Vaughan, \emph{Multiplicative Number Theory I},
Theorem 13.7 and Corollary 13.8,
\url{https://personal.science.psu.edu/rcv4/personal/Publications/MNTI/17.0_pp_419_451_Conditional_estimates.pdf}.
Only characters induced by conductors dividing the power-of-two
modulus are used, so the stated restricted GRH hypothesis suffices.

For some absolute $D>0$, this implies the following elementary
short-interval consequence: whenever $x$ is sufficiently large,
$Q\ge4$ is a power of two, and
\[
 H=DQ\sqrt x(\log x)^2\le x/2,
\]
every reduced class modulo $Q$ contains a prime power in $(x-H,x]$.
Indeed subtract the two $\psi$ estimates. The main term is
$H/\varphi(Q)=2D\sqrt x(\log x)^2$, whereas both error terms
together are bounded by a fixed multiple of
$\sqrt x(\log x)^2$. Choose $D$ larger than that fixed multiple.
Positive $\psi$ difference supplies a prime power. An actual prime
is unnecessary for the Paley construction, so no prime-power
subtraction estimate is needed.

\subsection{Construction and deficit estimate}
\begin{proof}[Proof of the theorem]
Fix a large multiple of four $n$. Set $Q=2^{T(n)}$ and
\[
 x=n/2-1,\qquad a\equiv n-1\pmod Q,\qquad
 H=DQ\sqrt x(\log x)^2.
\]
The residue $a$ is odd, hence coprime to $Q$. Since $Q=n^{o(1)}$,
we eventually have $Q\le\sqrt{x-H}$ and $H\le x/2$.
The conditional interval estimate gives a prime power
\[
 x-H<\ell\le x,\qquad \ell\equiv n-1\pmod Q.
\]
Because $4\mid Q$ and $4\mid n$, necessarily $\ell\equiv3\pmod4$.
Paley's type-I construction therefore gives a Hadamard matrix of
order $b=\ell+1$. Also
\[
 k=\frac{n-\ell-1}{Q}\in\mathbb Z_{>0},\qquad k\le n.
\]
The uniform-buffer lemma gives a Hadamard matrix of order
$a_1=Qk=n-b$. Here $b\le n/2$ and $a_1\ge n/2$.
Take the first $b$ rows of each matrix and concatenate horizontally.
Their row Gram matrices add to
$a_1I_b+bI_b=nI_b$, producing a partial Hadamard matrix with
\[
 r(n)\ge b>n/2-H.
\]
Finally
\[
 H\ll Q\sqrt n(\log n)^2
   =\sqrt n\exp\!\left(O\!\left(\frac{\log n}{\log\log n}\right)\right),
\]
since $\log\log n=o(\log n/\log\log n)$. This proves the claim.
\end{proof}

The progression residue is $n-1\pmod Q$, which need not equal
$-1\pmod Q$: the target order need not be divisible by the buffer.
This distinction is essential when $Q$ exceeds the power of two
dividing $n$.

More quantitatively, if
$\tau(u)\le A\log u/\log\log u+O(1)$, the same proof gives,
for every $\eta>0$ and all sufficiently large target orders,
\[
 r(n)\ge\frac n2-
 \sqrt n\exp\!\left((A\log2+\eta)
                 \frac{\log n}{\log\log n}\right).
\]
All fixed constants and logarithmic factors are absorbed by the
$\eta$ margin. No numerical threshold or explicit value of $A$
is claimed here.

The verifier supplies unconditional finite examples of the exact
concatenation step: it searches for actual primes in the required
residue, constructs both Hadamard blocks at known orders, and checks
the full rectangular Gram matrix as integers. These calculations
neither establish GRH nor verify the eventual prime-distribution
estimate or the asymptotic constants.
